% Round 20. AT20. family20. Candidate until exact independent review.
\section{Abelian-type strata, staircase products, and tetrahedral quartics}
\label{r20at:section}

\paragraph{Scope.}
This section gives a sufficient condition on actual Chow motives of
closed strata.  It is stronger than a numerical-motive hypothesis.
It is preserved by the specified staircase product construction, in
arbitrary relative dimensions.  Combined with known weight--monodromy
for curves and surfaces, it gives full-cycle G for any finite number
of the supplied surface factors and any supplied good coefficient.
An explicit simultaneous tetrahedral degeneration then gives dense
global tuples of quartic surfaces whose products have full heights.
The unrestricted product of arbitrary bad surfaces is not deduced
from this stronger stratum hypothesis.  The R9 product candidates
already contain related staircase and motivic ideas; no novelty or
historical-domain-exclusion claim is made here.

All schemes in the finite-field calculations are actual schemes,
all Chow groups and motives have rational coefficients, and every
degree is taken on a proper variety over its native finite field.

\subsection{The comparison prime for finitely many Chow motives}

Call a smooth projective $W/k$ of \emph{Chow abelian type} if its
rational Chow motive is an actual direct summand of a finite direct
sum of Tate twists of motives of abelian varieties over $k$.
Thus the inclusion and retraction are rational Chow correspondences
whose composite is the identity in Chow, not just numerically.
The class is closed under finite sums, summands, products, and
Tate twists.  We may split a finite number of constant fields by
one finite extension of $k$ before using this definition.

\begin{lemma}[One comparison prime for the whole finite diagram]
\label{r20at:commonprime}
Let $W_1,\ldots,W_t/k$ be of Chow abelian type.  There is a
prime $\ell\ne\operatorname{char}k$ such that every native
numerically zero rational cycle on every $W_i\times W_j$ has
zero geometric $\ell$-adic class, in every codimension.
\end{lemma}

\begin{proof}
Choose the actual motive presentations.  List all abelian carriers
$B_\alpha$ occurring in them and put
$B=\prod_\alpha B_\alpha^2$.  Every pair carrier
$B_\alpha\times B_\beta$ is a Chow summand of $B$, using the
zero sections and the projections on the unused factors; two copies
also cover $\alpha=\beta$.  Thus each $h(W_i\times W_j)$ is
a summand of a finite direct sum of twists of $h(B)$.

Clozel's Theorem~1 supplies a positive-density set of primes for
which numerical and geometric homological equivalence coincide
for all cycles on $B_{\bar k}$; choose one invertible prime from
that single set.  There is no intersection of unrelated
positive-density prime sets.  A native numerically zero class
remains numerically zero after every finite constant extension:
against a new complementary cycle use proper norm and projection,
with the full constant-field degree.  Consequently it is
geometrically numerically zero.

Apply the chosen Chow inclusion to a numerical-zero cycle on
$W_i\times W_j$.  Each component is an actual numerical-zero
cycle on a twist of $B$, hence has zero class at this prime.
Apply the Chow retraction.  Its composite with the inclusion is
the actual identity, so the original geometric class is zero.
All operations preserve numerical equivalence by the Chow
projection formula.  The statement uses neither full Tate
surjectivity nor a numerical-class-to-cohomology map.
\end{proof}

\begin{proposition}[An incidence certificate from component motives]
\label{r20at:inverse}
Let $\mathcal Y/R$ be regular projective reduced strictly
semistable, of relative dimension $d$, with finite residue field.
Suppose its components have Chow abelian type after one finite
constant extension.  After the corresponding unramified extension
there is an actual rational correspondence matrix
\[
 G\in\bigoplus_{i,j}\operatorname{CH}^{d-1}(D_j\times D_i)
 \quad\hbox{such that}\quad A^H G^H A^H=A^H
 \tag{AT20.1}\label{r20at:AGA}
\]
on full geometric cohomology at one invertible prime.
The conclusion is unchanged after tensoring the entire component
diagram with the diagonal of any smooth projective coefficient
$T_0/k$, without a comparison hypothesis on $T_0$.
\end{proposition}

\begin{proof}
PI20 gives an actual rational $G$ whose error $E=A-AGA$ is
numerically zero as a correspondence, using numerical-motive
semisimplicity.  Each entry of $E$ is an actual cycle on
$D_j\times D_i$.  Lemma~\ref{r20at:commonprime} kills all these
entries geometrically at the same prime, proving
\eqref{r20at:AGA}.  Tensoring gives the error
$\Delta_{T_0}\boxtimes E$, which also has zero geometric class.
Its action vanishes on every Kunneth summand, including mixed
classes.  PI20's corrected actual kernel lifts therefore give
$\Delta_A=\mathscr R_A=0$ in every degree on this product
diagram.  The good coefficient need not have a Chow decomposition.
\end{proof}

\subsection{Actual stratum blowups preserve the motive hypothesis}

Let $W/k$ be smooth projective and let $\partial W=\bigcup D_i$
be an SNC divisor.  By a closed stratum we mean a connected
component of some $D_J=\bigcap_{j\in J}D_j$, including $W$
for $J=\varnothing$.  Suppose all of these have Chow abelian type.

\begin{lemma}[All strata after a boundary-stratum blowup]
\label{r20at:stratumblowup}
Blow up the smooth intersection $Z=D_K$ of codimension $c=|K|$,
allowing its disjoint connected components simultaneously.  The reduced total
boundary on $\widetilde W=\operatorname{Bl}_Z W$ is SNC, and
every closed stratum of this boundary has Chow abelian type.
\end{lemma}

\begin{proof}
For $c=1$ the blowup is the identity.  Assume $c\ge2$ and use
the convention that $\mathbf P^{\rm lines}(N)$ parametrizes
one-dimensional subspaces of $N$.  Let $F$ be the exceptional
divisor, with
\[
 F=\mathbf P_Z^{\rm lines}(N_{Z/W}),\qquad
 N_{Z/W}=\bigoplus_{k\in K}\mathcal O_W(D_k)|_Z.
 \tag{AT20.2}\label{r20at:normal}
\]
The normal decomposition is the dual of the actual conormal decomposition of
the transverse Cartier equations.

For a collection $J$ of old boundary indices with $K\not\subset J$,
the nonempty intersection of the corresponding strict-transform
divisors is
\[
 \widetilde D_J
       =\operatorname{Bl}_{D_{J\cup K}}D_J.
 \tag{AT20.3}\label{r20at:strictstratum}
\]
Here the center has codimension $c-|J\cap K|$ in $D_J$;
a Cartier-center blowup is interpreted as the identity.  If
$K\subset J$, the intersection of all those strict transforms
is empty.  Its intersection with $F$ when nonempty is
\[
 F\cap\bigcap_{j\in J}\widetilde D_j
  =\mathbf P_{D_{J\cup K}}^{\rm lines}
       \left(\bigoplus_{k\in K\setminus J}
                   \mathcal O_W(D_k)|_{D_{J\cup K}}\right).
 \tag{AT20.4}\label{r20at:exceptionalstratum}
\]
These are scheme-theoretic formulas.  In the Rees charts they say
that each $j\in J\cap K$ sets one homogeneous normal coordinate
to zero, while each $j\notin K$ cuts the base transversely.
They also prove that the displayed new boundary is SNC.  If an
old intersection is disconnected, use its disjoint components;
the same equations apply to each.

The actual Chow blowup and projective-bundle formulas are
\[
 \begin{split}
 \operatorname{CH}^q(\operatorname{Bl}_Z W)
   &\simeq\operatorname{CH}^q(W)
       \oplus\bigoplus_{a=1}^{c-1}\operatorname{CH}^{q-a}(Z),\\
 (\alpha,\beta_a)&\longmapsto
 f^*\alpha+\sum_{a=1}^{c-1}
        j_*\bigl(\xi^{a-1}p^*\beta_a\bigr),\\
 \operatorname{CH}^q(\mathbf P^{\rm lines}(N))
   &\simeq\bigoplus_{a=0}^{\operatorname{rk}N-1}
                  \operatorname{CH}^{q-a}(Z),\quad
 (\beta_a)\longmapsto\sum_a\xi^a p^*\beta_a,
 \end{split}
 \tag{AT20.5}\label{r20at:chowformula}
\]
where $\xi=c_1(\mathcal O(1))$, $j:F\hookrightarrow\widetilde W$
and $p:F\to Z$.  They hold after product with every smooth
projective test variety.  Their maps and inverses are actual
Chow correspondences, so the identity principle gives the usual
direct-sum formulas in rational Chow motives.  Applying these to
\eqref{r20at:strictstratum}--\eqref{r20at:exceptionalstratum}
proves the claim for every new stratum.  This argument establishes
actual motive splittings, not merely cohomological decompositions.
\end{proof}

\subsection{A staircase component is an actual stratum-blowup tower}

\begin{proposition}[Closure for any finite number of supplied factors]
\label{r20at:productstrata}
Let $\mathcal Y_1,\ldots,\mathcal Y_m/R$ be regular projective
reduced strictly semistable models, with smooth geometrically
integral generic fibers.  Suppose every closed special-fiber
stratum of every factor has Chow abelian type after one finite
constant extension.  Their ordered staircase product has a
regular projective reduced strictly semistable model whose every
closed stratum has the same property.
\end{proposition}

\begin{proof}
It suffices to construct one binary step and verify the full
boundary, so that the step can be repeated.  Order the components
of the two factors globally.  At a point with $r$ and $s$ branches,
the product cone is
\[
 \{(a,b)\in\mathbf R_{\ge0}^{r+s}:\sum a_i=\sum b_j\},
 \qquad N=\{(a,b)\in\mathbf Z^{r+s}:\sum a_i=\sum b_j\}.
 \tag{AT20.6}\label{r20at:cone}
\]
Its height-one rays are $n_{ij}=(e_i,f_j)$.  Triangulate the
height-one product of simplices by the monotone paths in the
product order.  A path's first ray and successive differences
form an integral basis of $N$, so every resulting cone is
unimodular and every vertical multiplicity remains one.

This subdivision is projective.  Give vertex $(i,j)$ the height
$-ij$.  Transferring equal small mass from incomparable pairs
$(a,b'),(a',b)$ to $(a,b),(a',b')$, with $a<a'$ and $b<b'$,
preserves the two marginals and strictly decreases height.
A minimizing combination therefore has chain support.  The
monotone coupling gives every path simplex, and the heights
restrict compatibly on faces.  Thus they define a global coherent
subdivision for the fixed component orders.  Its monomial blowup
charts are those of Hartl's product resolution.  This yields an
actual projective model, not a merely formal local modification.

Fix its component continuing $D_i\times E_j$.  The quotient
star at $n_{ij}$ has the integral coordinates
\[
 n_{aj}\mapsto u_a,\qquad n_{ib}\mapsto v_b,\qquad
 n_{ab}\mapsto u_a+v_b.
 \tag{AT20.7}\label{r20at:star}
\]
Only the pairs $a<i,b<j$ or $a>i,b>j$ occur in its mixed rays.
In a positive orthant with basis $u_1,\ldots,u_M,v_1,\ldots,v_N$,
star-subdivide $\langle u_a,v_b\rangle$ by $u_a+v_b$, in
reverse product order.  Its maximal cones are the monotone
paths.  This follows by induction on $M+N$: the first ray
$u_M+v_N$ gives two cones omitting respectively $v_N$ and $u_M$;
the remaining steps restrict to the $(M,N-1)$ and $(M-1,N)$
problems.  The two-face for an unprocessed pair remains a face
until its own step, so each operation is a blowup of a smooth
codimension-two boundary stratum.

A path through $(i,j)$ splits into lower and upper paths with
disjoint coordinate sets.  Consequently the component map to
$D_i\times E_j$ is the actual sequence of blowups along the
proper transforms of
\[
           (D_i\cap D_a)\times(E_j\cap E_b),
 \tag{AT20.8}\label{r20at:centers}
\]
in ascending product order for the lower pairs and descending
product order for the upper pairs.  Globally use the same component
orders and the connected smooth pieces of these intersections.
Restricting this sequence to any chart gives exactly the quotient
star just computed.  The local identifications are the identity
on the common dense toroidal open, so glue over $D_i\times E_j$.
This argument has no bound on $r,s$ or on the relative dimensions.

The original pair on $D_i\times E_j$ has boundary given by all
old double divisors in either factor.  Its closed strata are
products of old closed strata.  Lemma~\ref{r20at:stratumblowup}
applies at each step of \eqref{r20at:centers}.
Finally its complete reduced boundary is exactly its intersection
with the other components of the new total model.  In the quotient
star every boundary ray is a two-dimensional cone containing
$n_{ij}$ and one other global component ray: the old axes and
the mixed rays in \eqref{r20at:star} give precisely these rays.
There is no further boundary ray or new vertical component.
Thus every actual closed stratum of the new total model is a
closed stratum of one of the just-treated component pairs.
This proves the binary assertion and allows induction on $m$.
No factorization theorem for arbitrary birational maps is used.
\end{proof}

\subsection{Full local cycle groups and all finite extensions}

\begin{theorem}[A good coefficient and supplied abelian-type surface strata]
\label{r20at:local}
Fix the finite valued $K/\mathbf Q_p$ and its absolute value.
Suppose $T/K$ is smooth, projective and geometrically integral,
has a smooth projective model, and
$S_1,\ldots,S_m/K$ have regular projective reduced strictly
semistable models, each $S_i$ being a smooth projective
geometrically integral curve or surface.  Suppose every closed
stratum of each of these factor models has Chow abelian type
after finite constant extension.  Then every geometrically
homologically trivial rational cycle on
\[
                         X=T\times\prod_{i=1}^mS_i
 \tag{AT20.9}\label{r20at:genericproduct}
\]
has strong ordinary-actual G in every $1\le r\le\dim X$.
The assertion holds for every new cycle after every finite valued
extension of $K$.  It also descends to an original valued field
if all of the displayed models are supplied over one common
finite extension of that field.
\end{theorem}

\begin{proof}
Pass first to one finite unramified extension defining the finitely
many motive splittings and splitting the stratum constants.
Proposition~\ref{r20at:productstrata} supplies the semistable
product witness for $P=\prod_iS_i$, with Chow-abelian components.
Apply Proposition~\ref{r20at:inverse}, then tensor its actual
inverse with the good coefficient $T_0$.  The resulting product
model for $X$ has $\Delta_A=\mathscr R_A=0$ in every degree
at one common prime.

All degrees of cohomology of each curve or surface satisfy
weight--monodromy.  The good factor is pure and unramified.
Kunneth and tensor compatibility of monodromy filtrations give
weight--monodromy on every degree of $X$, at every invertible
prime.  Thus OW19 kills the actual inertia image at the chosen
prime.  AK19/CR19 and FD17/VT17 give the stated G, with exact
specified support and all compact actual tests.  If $\dim X\le2$,
use the frozen divisor and zero-cycle endpoint theorems instead
of AK19's stated geometric range.

The same proof applies over every finite unramified extension of
this field: base change the actual Chow splittings and choose one
comparison prime anew.  Numerical-equivalence comparison is
applied to all cycles over that new field.  The product witness
is native and strictly semistable.  The frozen RI18 norm-detection
corollary therefore extends the full cycle assertion to every
finite valued extension, including ramified extensions and their
new cycles.  Concretely, in a finite Galois closure first apply
the theorem over the maximal unramified subfield and then RI18
to the totally ramified Galois step; descend to the desired
field by finite actual trace.  This uses the full extension
degrees, including wild ramification, and the extending original
absolute value.  No new cycle is assumed to be an old pullback.

If the hypotheses were supplied only over a common finite
extension $L/K_0$, a new cycle over any $K'/K_0$ is pulled to a
finite compositum with $L$, treated by the all-extension assertion
just proved, and descended with the complete field degree.  The
final cycle is the originally specified one.  No arbitrary proper
actual pushforward or finite-on-complement alteration is used.
\end{proof}

For a surface model the stratum hypothesis is automatic if all
its geometric components are rational or geometrically ruled:
double strata are curves and triple strata are points.  After
finite constant extension a ruled component is birational to a
projective bundle over a smooth projective curve.  Resolution of
a birational map of smooth projective surfaces by point blowups
and \eqref{r20at:chowformula} show that its motive is a summand
of curve motives and Tate twists.  This includes supplied models
of the usual ruled or rational Kulikov types.  Existence of such
a projective model is an explicit hypothesis; it is not asserted
for every K3 surface over every local field.

\subsection{An explicit tetrahedral surface with true triple strata}

\begin{proposition}[A quartic congruence with rational components]
\label{r20at:quartic}
Let the residue characteristic be odd and let $\pi$ be the
uniformizer.  Let $F$ be a quartic form over $R$ such that
$\bar F$ is nonzero at the four coordinate vertices of
$\mathbf P^3_k$ and cuts each of the six coordinate lines in
four simple geometric points away from their endpoints.  Then
\[
 \mathcal Q:\quad x_0x_1x_2x_3+\pi F(x_0,x_1,x_2,x_3)=0
                    \ \subset\mathbf P^3_R
 \tag{AT20.10}\label{r20at:quarticequation}
\]
has a projective regular reduced strictly semistable modification,
isomorphic on its generic fiber, with four geometrically rational
surface components, six double curves isomorphic to
$\mathbf P^1$, and four unchanged triple points.  Its generic
fiber is a smooth geometrically integral quartic K3 surface.
The conditions are satisfied by $\bar F=\sum_i x_i^4$.
\end{proposition}

\begin{proof}
The closed fiber consists of the four coordinate planes.  At a
point on just one plane the total space is smooth over $R$.
At a double-line point where $\bar F\ne0$, dividing by units
gives a regular semistable equation $xy=\pi u$.  At a coordinate
vertex $\bar F$ is a unit, and the equation has the regular
semistable form $xyz=\pi u$ with three branches.

The only other points are the $24$ simple geometric zeros of
$\bar F$ on the double lines.  Write the local equation as
$u xy+\pi F=0$ with $u$ invertible.  The simple-zero assumption
makes $z=F/u$ a local etale coordinate along the line.  The actual
etale local equation is therefore
\[
                              xy+\pi z=0.
 \tag{AT20.11}\label{r20at:node}
\]
Blow up the ideal $(x,\pi)$ of one of its two plane components.
The actual Rees charts, with the factors removed by the relation
inside the common fraction field, are
\[
 \begin{array}{c|c|c}
 \text{chart}&\text{coordinates and relations}&\text{map}\\ \hline
 x=\pi t& R[t,y],&x=\pi t,\ z=-ty,\\
 \pi=xv& R[x,v,z]/(xv-\pi),&y=-vz.
 \end{array}
 \tag{AT20.12}\label{r20at:smallresolution}
\]
The first chart is smooth over $R$ and the second is regular
strictly semistable.  The chosen plane is blown up at the point,
the other plane is unchanged, and their double curve remains
isomorphic to the original line.  The exceptional curve is contained
in the chosen vertical surface; it is not an additional vertical
surface component.

Globally blow up the ideals of the coordinate-plane components in
the order $0,1,2$, always using their proper transforms.  Away
from the remaining points \eqref{r20at:node}, each such component
is Cartier, including at the triple points, so its blowup is an
isomorphism there.  At a remaining node the displayed charts
resolve it.  Every double line has an endpoint component among
$0,1,2$, so every node is resolved.  The four final components
are the planes blown up at respectively $12,8,4,0$ geometric
points.  Closed points of larger residue degree are blown up
natively, and split after a finite unramified extension.
The double curves remain $\mathbf P^1$ and the four vertices
are unchanged.  This proves the asserted regular reduced
semistable model with actual triple strata.

The semistable charts cover the closed fiber.  Since the model
is proper, their union covers the entire model: its closed
complement would have proper image contained in the generic
point of the trait, which is impossible.  Its generic fiber is
therefore smooth.  A smooth quartic surface in $\mathbf P^3$
is geometrically integral and is K3 by adjunction and the
standard hypersurface cohomology sequence.  Finally, for odd
residue characteristic $x^4+y^4$ has four simple geometric
zeros with neither coordinate zero, proving the asserted
nonempty coefficient condition.
\end{proof}

\paragraph{The generic fiber is not potentially good.}
The dual complex of the supplied model is the boundary of a
tetrahedron, so its second cohomology is $\mathbf Q_\ell$.
The weight spectral sequence identifies this with the nonzero
weight-zero piece of generic $H^2$.  Surface weight--monodromy
then gives $N^2\ne0$.  A finite extension multiplies $N$ by
its ramification index and cannot kill it, whereas good reduction
has $N=0$.  For a product of $m$ such surfaces, on the tensor
of their $H^2$ spaces one has
\[
 N^{2m}(v_1\otimes\cdots\otimes v_m)
   =\frac{(2m)!}{(2!)^m}
       N_1^2v_1\otimes\cdots\otimes N_m^2v_m\ne0
 \tag{AT20.13}\label{r20at:monodromy}
\]
after choosing $N_i^2v_i\ne0$.  Thus this is not a change of
model of a potentially good generic fiber.

\subsection{Dense global tuples and full complementary heights}

\begin{theorem}[Dense quartic products with a common deep bad place]
\label{r20at:global}
Fix an odd rational prime $p$ and an integer $m\ge1$.
There is a Zariski-dense set of ordered $m$-tuples of quartic
forms over $\mathbf Q$ defining smooth surfaces
$S_1,\ldots,S_m$ such that $X=\prod_iS_i$ has full
strong G at every finite completion and after every finite
number-field extension, for all new homological rational cycles.
Consequently PH16 has the complete complementary domains
\[
 \operatorname{CH}^r(X_E)_{\mathbf Q,\mathrm{geom\,hom}}
 \times
 \operatorname{CH}^{2m+1-r}(X_E)_{\mathbf Q,\mathrm{geom\,hom}}
                          \longrightarrow\mathbf R
 \tag{AT20.14}\label{r20at:height}
\]
for every number field $E/\mathbf Q$ and $1\le r\le2m$.
All factors have the tetrahedral degeneration at $p$, and
$N^{2m}\ne0$ for the product there.
\end{theorem}

\begin{proof}
Work in the affine space of the $35$ quartic coefficients.  Choose
the first form, and then each next form, in the integral congruence
class
\[
 P_i\equiv x_0x_1x_2x_3+p\sum_jx_j^4\pmod {p^2}.
 \tag{AT20.15}\label{r20at:pcongruence}
\]
After earlier choices, let $B$ be the finite union, excluding $p$,
of primes where their chosen primitive projective quartic closures
are not smooth.  At every $q\in B$ impose that the next form
reduce to a fixed geometrically smooth quartic over $\mathbf F_q$.
Such a form exists for every $q$: use the Fermat quartic for odd
$q$, and for $q=2$ use
\[
        H=x^4+x^3y+y^3z+z^3w+w^3x.
 \tag{AT20.16}\label{r20at:char2good}
\]
Indeed, if its four partial derivatives vanish, a zero coordinate
would force all coordinates zero.  Otherwise the four cyclic
terms $x^3y,y^3z,z^3w,w^3x$ are equal, so sum to zero in
characteristic two, while $x^4\ne0$.  Such a critical point
does not lie on $H=0$.  This proves geometric smoothness.

The Chinese remainder theorem makes all these coefficient
conditions compatible and gives a full-rank integral affine
progression.  It is Zariski dense; choose any point in it
defining a smooth generic quartic.  In fact the $p$-condition
already gives smoothness by Proposition~\ref{r20at:quartic}.
The finite bad set exists because a smooth projective generic
hypersurface spreads to a smooth projective model outside
finitely many primes.  At a prime $q\ne p$ at most one
factor is bad: any later factor was required to be good at
every bad prime of the earlier ones.
The content of a chosen integral form is prime to $p$ by
\eqref{r20at:pcongruence}; dividing by it changes no local
projective model at $p$.  At a specified good prime its content
is also a unit, since the prescribed smooth reduction is nonzero.

At $p$, all factors have the model of
Proposition~\ref{r20at:quartic}, whose strata have Chow abelian
type after splitting the finitely many closed points.  Apply
Theorem~\ref{r20at:local}, with the coefficient a point, to
all factors simultaneously.  Its RI18 conclusion treats every
finite extension of $\mathbf Q_p$.  At any $q\ne p$, the
product of the good factors has a smooth projective model;
if one surface is bad, AS20 applies with that arbitrary
surface as the bad factor.  If none is bad, use good reduction
or the same AS20 statement.  These arguments remain valid
over every finite completion of every number field and for
all new cycles.  PH16 and GA17 now give
\eqref{r20at:height} in every complementary degree.

For density of the tuple set, use induction on $m$.  For each
already chosen $(m-1)$-tuple, the permitted next forms contain
a full-rank affine progression, hence are Zariski dense in
the last coefficient space.  A polynomial vanishing on all
constructed tuples therefore has every last-variable
coefficient vanishing on all earlier constructed tuples.
The inductive density makes every coefficient zero.
The assertion is density of these explicitly constrained
tuples, not a theorem about all quartic surfaces or all tuples.
The monodromy assertion is \eqref{r20at:monodromy}.
\end{proof}

\paragraph{Dependencies, history and remaining scope.}
Numerical semisimplicity and the finite error are PI20, based
on Jannsen's Theorem~1.  The comparison-prime input is Clozel,
\emph{Equivalence numerique et equivalence cohomologique pour
les varietes abeliennes sur les corps finis}, Ann. of Math.
150 (1999), 151--163, Theorem~1,
\url{https://arxiv.org/pdf/math/9907203}; it is used for one
finite carrier product.  The model input is Hartl,
\emph{Semi-stability and base change}, Arch. Math.
77 (2001), 215--221, Proposition~2.1 and its actual monomial
charts, \url{https://ivv5hpp.uni-muenster.de/u/uhartl/Publikat/SemistBC.pdf}.
The fan calculation is written above and also appears, for
surface products, in the retained R9 ST candidate and the
current PX20 investigation.  The actual Chow formulas are
the classical projective-bundle and smooth-blowup formulas;
their use after every test product gives the motive statement.
Surface weight--monodromy is the known input recorded in OW19
and Saito's discussion after Conjecture~0.5.  Arbitrary finite
new cycles use the exact frozen RI18 norm theorem, rather
than an unproved preservation of motives under arbitrary
ramified resolutions.

The motive condition is a sufficient supplied-model hypothesis.
It is not deduced for components of arbitrary semistable
alterations, and it is not required for general G.  Nothing
here proves the general incidence error homologically zero,
or proves general G or full heights on an arbitrary variety.
The global family is a concrete consequence under verified
local degeneration conditions; no external first-discovery
claim is made.
